\section{The Empirical Shift: Falling Inference Cost and Deployment Migration}

The relevant empirical question is not whether AI capability has stopped improving---it has not---but whether the economic relationship between frontier capability and deployment-relevant capability is changing. Three types of evidence must be distinguished: measured traffic, inferred deployment patterns, and claims about market power. The first is direct; the second is an interpretation of the first; the third goes beyond what the data can establish on its own.

\textbf{Measured traffic: the cost of fixed capability is collapsing.} Epoch AI finds that the cost of achieving a given capability level has fallen roughly 47\% per quarter---about 13x per year---since 2023, with the decline steepest immediately after a capability becomes state-of-the-art (66\% per quarter initially, decaying to 32\% two years later)~\cite{epochai2026}. The premium for frontier capability is therefore short-lived: any given capability is rapidly commoditized, and the frontier lab's advantage is a moving target that must be continuously re-earned at rising cost. This is a measured trend in inference pricing; market power is a separate claim.

\textbf{Measured traffic: deployment is bifurcating---in gateway traffic, with reasonable inference about deployment at large.} Vercel's AI Gateway Production Index, based on traffic through September 2026, reports that open-weight models accounted for 71.5\% of gateway token volume and approximately 25\% of estimated spend, up from roughly 5\% of spend in June~\cite{vercel2026}. The September index, reporting August data, showed 56\% of tokens and 14\% of spend---a one-month jump of 15.5 percentage points in token share and 9 points in spend share. DeepSeek-v4.1-Flash alone processed nearly half of September's tokens while accounting for roughly 2.6\% of gateway spend. Yet 68\% of teams that increased spending on open-weight models also increased spending on closed-weight models, and frontier spend grew approximately 13-fold since December while average estimated cost per token fell 47\%; nearly one-third of teams changed their most-used lab in September. The pattern is a split, not a replacement: high-volume, low-stakes workloads migrate to commoditized open models; high-value, high-reliability workloads remain with frontier labs, which grow in absolute dollars but lose relative share. This article uses the rounded spend figure of ``approximately 25\%'' to match the source's own framing (``nearly 25\%'' in the Vercel blog post); an earlier draft used 23\%, which slightly understated the reported value.

\textbf{Inferred deployment patterns and the limits of the data.} The figures describe Vercel AI Gateway traffic, not the entire AI market, and rely on list-price estimates and a model classification that has broadened over time; they are used here as a convergent indicator, not a census. Token volume does not establish the economic value of tasks, model quality, or the profitability of each segment. The assertion that frontier labs lose pricing power is plausible but requires separate evidence; absolute spending growth and falling average token costs do not establish it on their own. The September-to-October comparison carries a classification caveat: the open-weight definition broadened over time, so the June-to-September series is not perfectly consistent. These limitations bound what can be claimed: the data show migration of gateway traffic volume and estimated spend, which is consistent with (but does not independently establish) deployment migration at the market level.

\textbf{The mechanism is C$\rightarrow$B.} Falling inference cost for fixed capability makes free-riding on open models rational for the commoditized layer; the deployment shift is the observable consequence. As the economic gradient steepens---frontier costs rise faster than frontier returns---the rational follower strategy shifts from ``catch the frontier'' to ``deploy the commoditized capability,'' and the capability gradient closes from below, jaggedly, not uniformly. Open-weight models lagged the closed frontier by roughly four months on Epoch's benchmark index as of mid-2026, a gap that has not consistently narrowed~\cite{epochai2026}. This is expected: the mechanism predicts deployment migration, not benchmark convergence.

\textbf{Falsifiability.} The mechanism predicts that if the ratio of open-weight spend share to token share stopped narrowing---if open-weight models captured volume without capturing value---the premise that follower free-riding is economically rational would weaken, and the rentier turn would not be predicted. The observed narrowing (from roughly 4:1 in August to roughly 3:1 in September) is consistent with the mechanism; a reversal would challenge it. This is a claim-level falsification condition, not a test of the entire political-economy argument.

\textbf{What this does not establish.} Strong benchmark results support continued progress in particular capabilities; they do not independently establish a Perez-style transition. Improvements in benchmark performance can coexist with rising training costs, expensive inference, and persistent failures in real-world settings. Conversely, diminishing returns from increasing model size do not establish that progress as a whole is slowing. Additional inference-time computation, architectural changes, synthetic data, tool use, and system-level engineering may change the relationship between resources and performance. The analytically important variable is not capability in isolation, but the cost and reliability of converting additional resources into useful capability.
